% !TeX program = xelatex
\documentclass[a4paper,10pt]{article}

% ============================================================================
% ក្រដាសដាច់ដោយឡែក៖ សេណារីយ៉ូនៃការបង្រៀន (Scenario នៃការបង្រៀន – ៥០ នាទី)
% ============================================================================
\usepackage[a4paper, left=0.95cm, right=0.95cm, top=1.0cm, bottom=1.1cm, footskip=0.5cm]{geometry}
\usepackage{amsmath}
\usepackage{unicode-math}
\usepackage{fontspec}
\usepackage{array, longtable, colortbl}
\usepackage[table]{xcolor}
\usepackage{pifont}
\usepackage{fancyhdr}
\usepackage{tikz}

% LaTeX 2025 + longtable compatibility guard
\makeatletter
\let\@expl@@@mark@update@singlecol@structures@@\relax
\makeatother

% Cross-platform font path resolver (Windows local + Linux/Cloud repo fonts/)
\providecommand{\NGSFindFonts}{%
  \IfFileExists{C:/Windows/Fonts/cambria.ttc}{%
    \def\NGSSysPath{C:/Windows/Fonts/}%
    \def\NGSUserPath{C:/Users/KIT/AppData/Local/Microsoft/Windows/Fonts/}%
  }{%
    \IfFileExists{./fonts/KhmerOSsiemreap.ttf}{\def\NGSSysPath{./fonts/}\def\NGSUserPath{./fonts/}}{%
    \IfFileExists{../fonts/KhmerOSsiemreap.ttf}{\def\NGSSysPath{../fonts/}\def\NGSUserPath{../fonts/}}{%
    \IfFileExists{../../fonts/KhmerOSsiemreap.ttf}{\def\NGSSysPath{../../fonts/}\def\NGSUserPath{../../fonts/}}{%
    \IfFileExists{../../../fonts/KhmerOSsiemreap.ttf}{\def\NGSSysPath{../../../fonts/}\def\NGSUserPath{../../../fonts/}}{%
    \IfFileExists{../../../../fonts/KhmerOSsiemreap.ttf}{\def\NGSSysPath{../../../../fonts/}\def\NGSUserPath{../../../../fonts/}}{%
    \IfFileExists{../../../../../fonts/KhmerOSsiemreap.ttf}{\def\NGSSysPath{../../../../../fonts/}\def\NGSUserPath{../../../../../fonts/}}{%
    \IfFileExists{../../../../../../fonts/KhmerOSsiemreap.ttf}{\def\NGSSysPath{../../../../../../fonts/}\def\NGSUserPath{../../../../../../fonts/}}{%
      \def\NGSSysPath{../../../../../../../fonts/}\def\NGSUserPath{../../../../../../../fonts/}%
    }}}}}}}%
  }%
}
\NGSFindFonts

\setmathfont[
  Path={\NGSSysPath},
  FontIndex=1,
  Scale=1.02
]{cambria.ttc}

\defaultfontfeatures{Ligatures=TeX}

\setmainfont[
  Path={\NGSUserPath},
  Extension=.ttf,
  Script=Khmer,
  BoldFont={KhmerOSsiemreap},
  BoldFeatures={FakeBold=1.8},
  ItalicFont={KhmerOSsiemreap},
  ItalicFeatures={FakeSlant=0.25},
  BoldItalicFont={KhmerOSsiemreap},
  BoldItalicFeatures={FakeBold=1.8, FakeSlant=0.25}
]{KhmerOSsiemreap}

\newfontfamily\muol[
  Path={\NGSUserPath},
  Extension=.ttf,
  Script=Khmer
]{KhmerOSmuollight}

\newfontfamily\TimesEn[
  Path={\NGSSysPath},
  Extension=.ttf,
  Scale=1.02,
  BoldFont={timesbd},
  ItalicFont={timesi},
  BoldItalicFont={timesbi}
]{times}

\newfontfamily\NagaBold[
  Path={\NGSUserPath},
  Extension=.ttf
]{AKbalthom Naga Bold}

\newfontfamily\CalibriFont[
  Path={\NGSSysPath},
  Extension=.ttf
]{calibri}

\XeTeXlinebreaklocale "km"
\XeTeXlinebreakskip = 0pt plus 1pt minus 0.1pt

\definecolor{stepgray}{HTML}{D9D9D9}
\definecolor{headgray}{HTML}{F2F2F2}
\definecolor{darkred}{HTML}{C00000}
\definecolor{red5e}{HTML}{FF0000}
\definecolor{figblue}{HTML}{1F4E9E}
\definecolor{tealink}{HTML}{0F766E}

\pagestyle{fancy}
\fancyhf{}
\renewcommand{\headrulewidth}{0pt}
\fancyfoot[R]{\CalibriFont\fontsize{11pt}{13pt}\selectfont\thepage}

\setlength{\parindent}{0pt}
\setlength{\parskip}{0pt}
\setlength{\tabcolsep}{4pt}

\newcolumntype{L}[1]{>{\raggedright\arraybackslash}p{#1}}
\newcolumntype{C}[1]{>{\centering\arraybackslash}p{#1}}
\newcommand{\en}[1]{{\TimesEn #1}}
\newcommand{\enbf}[1]{{\TimesEn\textbf{#1}}}
\newcommand{\fiveE}[1]{{\NagaBold\fontsize{9.5pt}{13pt}\selectfont\textcolor{red5e}{#1}}}
\newcommand{\ditem}{\tikz[baseline=-0.5ex]\fill (0,0.40ex) -- (0.40ex,0) -- (0,-0.40ex) -- (-0.40ex,0) -- cycle;\hspace{0.25em}}
\newcommand{\chkbox}{\tikz[baseline=-0.1ex]\draw[thin](0,0) rectangle (0.24,0.24);}
\newcommand{\cnum}[1]{{\fontsize{10pt}{11.5pt}\selectfont\ding{\numexpr181+#1\relax}}\hspace{0.15em}}

\begin{document}
\fontsize{9.8pt}{14.6pt}\selectfont

% -------------------- ក្បាលឯកសារ Scenario --------------------
\begin{center}
  {\muol\fontsize{11pt}{16pt}\selectfont ព្រះរាជាណាចក្រកម្ពុជា}\\[1pt]
  {\muol\fontsize{10pt}{15pt}\selectfont ជាតិ សាសនា ព្រះមហាក្សត្រ}\\[4pt]
  {\muol\fontsize{10.5pt}{15pt}\selectfont សេណារីយ៉ូនៃការបង្រៀន \enbf{(Scenario} នៃការបង្រៀន\enbf{)}}
\end{center}

\vspace{2pt}
{\renewcommand{\arraystretch}{1.08}
\begin{tabular}{@{}l l l p{1.0cm} l l l@{}}
  - អនុវិទ្យាល័យ & : & ព្រះក្រសាំង & & - មុខវិជ្ជា & : & គណិតវិទ្យា ថ្នាក់ទី៨ \\
  - មេរៀនទី១២ & : & ការប្រៀបធៀបត្រីកោណ (១.២.២ ករណី ជ.ម.ជ) & & - រយៈពេលសរុប & : & \textcolor{darkred}{\textbf{១ ម៉ោងសិក្សា (៥០ នាទី)}} \\
  - គ្រូបង្រៀន & : & ហុក ប៊ុនថែន & & - គំរូបង្រៀន & : & 5Es សកម្មភាពច្នៃប្រឌិត \\
\end{tabular}
}

\vspace{4pt}

% -------------------- I. តារាងបែងចែកពេលវេលា និងសេណារីយ៉ូលម្អិត (៥០ នាទី) --------------------
\textbf{I.} {\muol តារាងបែងចែកពេលវេលា និងសេណារីយ៉ូប្រតិបត្តិការក្នុងថ្នាក់ \enbf{(}៥០ នាទី\enbf{)}}

\vspace{2pt}
\renewcommand{\arraystretch}{1.14}
\begin{longtable}{|c|L{2.65cm}|c|L{7.15cm}|L{4.15cm}|}
\hline
\rowcolor{headgray}
\multicolumn{1}{|c|}{\muol នាទី} &
\multicolumn{1}{C{2.65cm}|}{{\muol ជំហាន \enbf{/ 5Es}}} &
\multicolumn{1}{c|}{\muol រយៈពេល} &
\multicolumn{1}{C{7.15cm}|}{{\muol សេណារីយ៉ូប្រតិបត្តិ \enbf{(}ពាក្យពេចន៍គ្រូ \enbf{\&} សម្ភារៈ\enbf{)}}} &
\multicolumn{1}{C{4.15cm}|}{{\muol ការគ្រប់គ្រង \enbf{\&} វិធីបម្រុង}} \\
\hline
\endfirsthead

\hline
\rowcolor{headgray}
\multicolumn{1}{|c|}{\muol នាទី} &
\multicolumn{1}{C{2.65cm}|}{{\muol ជំហាន \enbf{/ 5Es}}} &
\multicolumn{1}{c|}{\muol រយៈពេល} &
\multicolumn{1}{C{7.15cm}|}{{\muol សេណារីយ៉ូប្រតិបត្តិ \enbf{(}ពាក្យពេចន៍គ្រូ \enbf{\&} សម្ភារៈ\enbf{)}}} &
\multicolumn{1}{C{4.15cm}|}{{\muol ការគ្រប់គ្រង \enbf{\&} វិធីបម្រុង}} \\
\hline
\endhead

\hline
\endfoot

\hline
\endlastfoot

00--02 & \textbf{ជំហានទី១}\newline រដ្ឋបាលថ្នាក់ & ២ នាទី &
\ditem ហៅអវត្តមានរហ័ស ពិនិត្យអនាម័យតុ ។ &
ចាត់តាំងប្រធានក្រុមទុកមុន ដើម្បីកុំឱ្យខាតពេលរើតុកៅអី។ \\[3pt]
 &  &  &
\ditem ចែកសិស្សជាក្រុម ៣ នាក់ («ក្រុមវិស្វករ») ដាក់កញ្ចប់សម្ភារៈពិសោធន៍ទុកមុខតុ (ហាមទាន់បើកមុនបញ្ជា)។ &
ប្រកាសពិន្ទុវិន័យសម្រាប់ក្រុមគោរពច្បាប់។ \\ \hline

02--05 & \textbf{ជំហានទី២ (ក)}\newline កំណែកិច្ចការផ្ទះ & ៣ នាទី &
\ditem បញ្ចាំង\textbf{ល្បែងចាប់កំហុស} ៖ ដាក់សំណួរបញ្ឆោតលើ \en{Slide} «តើ $\angle KBD = \angle ABD$ ត្រូវឬខុស?» &
មិនចាំបាច់ហៅសិស្សសរសេរពេញលើក្ដារខៀនទេ។ \\[3pt]
 &  &  &
\ditem ទុកពេល ៣០ វិនាទីឱ្យដៃគូពិភាក្សា រួចសង្ខេបចម្លើយ\textbf{«ម.ជ.ម»} ទាំង ២ សំណួរ ក និង ខ លើ \en{Slide}។ &
ផ្ដោតតែត្រង់មុំបន្ថែមគ្នា $180^\circ$។ \\ \hline

05--08 & \textbf{ជំហានទី២ (ខ)}\newline រំលឹកមេរៀនទាក់ទង & ៣ នាទី &
\ditem \textbf{កាយវិការរហ័ស ១ (១.៥ នាទី)} ៖ បញ្ជា «លើកប៊ិច ២ ខ្វែងគ្នា!» សួរពីមុំទល់កំពូល + បញ្ចាំង \en{Slide} បញ្ជាក់ $\angle AOB = \angle A'OB'$ ។ &
បើសិស្សភ្លេចឈ្មោះ «មុំទល់កំពូល» ឱ្យធ្វើកាយវិការដៃខ្វែងម្ដងទៀតព្រមគ្នា។ \\[3pt]
 &  &  &
\ditem \textbf{កាយវិការរហ័ស ២ (១.៥ នាទី)} ៖ បញ្ជា «បត់កែងក្រដាសព្រៀងជាពីរ!» ភ្ជាប់ទៅកន្លះបន្ទាត់ពុះមុំ $OZ$ ។ &
ភ្ជាប់ទៅសញ្ញា \cnum{2} ក្នុងធនាគារហេតុផល។ \\ \hline

08--11 & \fiveE{E1 Engage}\newline បេសកកម្មស្រះទឹក & ៣ នាទី &
\ditem និទានបញ្ហាពិត៖ «វិស្វករចង់វាស់ទទឹងស្រះ $AB$ ដោយមិនចុះទឹក គាត់ដៅបង្គោល $O$ វាស់ $OA'=OA$, $OB'=OB$ និង $A'B'=18\text{ m}$ ។ តើគាត់ហ៊ានថា $AB=18\text{ m}$ ដោយសំអាងលើអ្វី?» &
កុំទាន់ផ្ដល់ចម្លើយពិតនៅទីនេះ ទុកឱ្យសិស្សចង់រកចម្លើយនៅ \en{E2--E5}។ \\[3pt]
 &  &  &
\ditem ស្រង់មតិសិស្សទស្សន៍ទាយ (ជឿ/មិនទាន់ជឿ) រួចសរសេរចំណងជើងមេរៀនថ្មី ១.២.២។ &
កត់សម្លេងមតិទុកលើគែមក្ដារខៀន។ \\ \hline

11--22 & \fiveE{E2 Explore}\newline ពិសោធន៍ ១ \& ២ & ១១ នាទី &
\ditem \textbf{ពិសោធន៍ ១៖ ចង្កឹះគន្លឹះ \& ក្រដាសស្តើង (៧ នាទី)} ៖ បើក–បិទ $\angle A$ ($30^\circ\leftrightarrow 100^\circ$) ពោលព្រមគ្នា «ជ្រុង $BC$ ចាក់សោពេលមុំអមចន្លោះ $\angle A$ ចាក់សោ»។ &
គ្រូដើរត្រួតការចាក់សោ $60^\circ$ និងការប្រើបន្ទះ \en{1:1} ផ្ទាល់ដៃ។ \\[3pt]
 &  &  &
+ ៥ នាទី៖ ចាក់សោ $\angle A=60^\circ$ គូស $BC \approx 5.29\text{ cm}$ រួចកាត់/គូរលើក្រដាសស្តើង \textbf{បង្វិល–ត្រឡប់ត្រួតបេះបិទ} លើត្រីកោណក្រុមជិតខាង ។ &
បើក្រុមយឺត ប្រើបន្ទះគំរូកាត់ស្រេចជួយផ្ទៀង។ \\[3pt]
 &  &  &
\ditem \textbf{ពិសោធន៍ ២៖ \en{POE} អន្ទាក់ ជ.ជ.ម (៤ នាទី)} ៖ ៣០វិ. ទស្សន៍ទាយ (\en{P}) $\rightarrow$ ២ន. យោលបន្ទះ $BC=4\text{ cm}$ ផ្ទាល់ដៃកាត់ $Ax$ ត្រង់ $C_1, C_2$ (\en{O}) $\rightarrow$ ១.៥ន. ពន្យល់ + បញ្ចាំង \en{Slide} សំយោគ (\en{E}) ។ &
ការពារការប្រើបន្ទះបញ្ឆោត $5\text{ cm}$ ច្រឡំមុនពេលកំណត់។ \\ \hline

22--28 & \fiveE{E3 Explain}\newline បង្គុំទ្រឹស្ដីបទ & ៦ នាទី &
\ditem បញ្ជាកាយវិការបើក–ចាក់សោ \en{S--A--S} នាំបង្គុំទ្រឹស្ដីបទ ជ.ម.ជ លើរូបបង្វិល–ឆ្លុះ $140^\circ$ ។ &
សង្កត់ធ្ងន់៖ កំពូលមុំត្រូវជាប់ចន្លោះជ្រុងទាំងពីរជានិច្ច ទោះរូបដាក់បញ្ច្រាស ឬត្រួតលើគ្នា។ \\[3pt]
 &  &  &
\ditem បញ្ជាក់ \textbf{គោលការណ៍កំពូលរួម} និងបិទ \textbf{ធនាគារហេតុផល} (\cnum{1} បម្រាប់ \cnum{2} មុំទល់កំពូល/ពុះមុំ \cnum{3} ជ្រុងរួម \cnum{4} ជ.ម.ជ $\rightarrow$ \cnum{5} ធាតុត្រូវគ្នា) ។ &
ចង្អុលទៅផ្ទាំងយោងប្រចាំថ្នាក់ (\en{Printable} ទំព័រ ៤)។ \\[3pt]
 &  &  &
+ ញែកសញ្ញា $=$ (ប្រវែង/រង្វាស់) vs $\cong$ (ត្រីកោណប៉ុនគ្នា) ។ &
ត្រួតពិនិត្យសិស្សមិនឱ្យសរសេរ $AB \cong A'B'$ ។ \\ \hline

28--38 & \fiveE{E4 Elaborate}\newline កាត ៣ ជំហាន (ស-ខ្មៅ) & ១០ នាទី &
\ditem ចែក \textbf{កាតស៊ើបអង្កេតសញ្ញាគន្លឹះ ៣ ជំហាន (ស-ខ្មៅ)} ៖ &
តម្រូវឱ្យម្រាមដៃចង្អុលអូសលើរូប ស-ខ្មៅ មុនគូសសញ្ញា។ \\[3pt]
 &  &  &
+ ៣ នាទី៖ ចង្អុលអូស + គូស $|$ ($OA=OB$), ចង្អុលបត់ + ព័ទ្ធកំពូល $O$ \& ធ្នូ ($\angle AOM=\angle BOM$), ចង្អុលអូស + គូស $\parallel$ ($OM=OM$) &
សង្កេតក្រុមដែលភ្លេចកំណត់កំពូល $O$ ។ \\[3pt]
 &  &  &
+ ៥ នាទី៖ ផែនការ \en{P2} (រក $\triangle$ ផ្ទុក $MA,MB$) + សរសេរ ៣ បន្ទាត់ដាក់លេខ \cnum{1}--\cnum{4} (បញ្ជាក់ \cnum{3} ជ្រុងរួម) + វិបាក \cnum{5} &
បើសិស្សភ្លេច $OM=OM$ ឱ្យម្រាមដៃចង្អុលអូសកំណាត់កណ្ដាល $OM$ (\cnum{3} \en{Reflexive})។ \\[3pt]
 &  &  &
+ ២ នាទី៖ \textbf{\en{POE} ៣០វិ.} $M$ រំកិលលើ $OZ$ (ញែកអ្វីប្រែ vs អ្វីថេរ) ។ &
ក្រុមលឿនបន្តធ្វើកាត ២--៤ លើសន្លឹកដដែល។ \\ \hline

38--48 & \fiveE{E5 Evaluate}\newline \en{Relay} \& មេដៃ ៤ កាត & ១០ នាទី &
\ditem \textbf{\en{Relay Proof Race} ៣ តួនាទី (៧ នាទី)} ៖\newline
+ អ្នកទី១ (\textbf{\en{Pause-Ask}})៖ ពន្យល់មាត់មុនគូស (ចាប់ \cnum{2} $\angle AOB=\angle A'OB'$, \textbf{ហាមគូស $AB=A'B'$}) ព័ទ្ធ $O$ &
ត្រៀមកាតពិន្ទុ ៣ ក្រុមដំបូង ការពារកំហុសគូស $AB=A'B'$ មុនស្រាយ។ \\[3pt]
 &  &  &
+ អ្នកទី២៖ ផ្ទៀង $O$ ចន្លោះជ្រុង សរសេរ \cnum{1}--\cnum{4} (រក្សា $=$ vs $\cong$)\newline
+ អ្នកទី៣ (\textbf{\en{Pause-Review}})៖ ផ្ទៀងលំដាប់ $O\to A\to B$ ទល់ $O\to A'\to B'$ ទាញ $AB=A'B'=18\text{ m}$ (\cnum{5}) បិទ \en{E1} ។ &
ផ្ទៀងលំដាប់កំពូលឱ្យត្រូវតាមជ្រុងប៉ុន។ \\[3pt]
 &  &  &
\ditem \textbf{មេដៃបញ្ឆោត ៤ កាត ស-ខ្មៅ (៣ នាទី)} ៖ ចង្អុលខ្យល់ \en{S--A--S} មុនលើកមេដៃ/ខ្វែងដៃ (ជ.ម.ជ $140^\circ$, ជ.ជ.ម, គ្មានសញ្ញាស្មានភ្នែក, ម.ម.ម/លំដាប់ខុស) + ម្រាម ១--៣ ទំនុកចិត្ត ។ &
បើសល់ពេលតិច បង្រួមមេដៃបញ្ឆោតមក ១.៥--២ នាទី (៤ កាតរហ័ស)។ \\ \hline

48--50 & \textbf{ជំហានទី៥}\newline កិច្ចការផ្ទះ \& បណ្ដាំផ្ញើ & ២ នាទី &
\ditem បញ្ចាំងកិច្ចការផ្ទះ ($\triangle ABC$ សមបាត $AB=AC$, បន្ទាត់ពុះ $AD$, $E\in AD$, ស្រាយតាមធនាគារ \cnum{1}--\cnum{5}) និងផ្ញើចូល \en{Telegram / Classroom} ។ &
កុំបន្លាយលើសម៉ោង បញ្ជូនរូបលំហាត់តាមអនឡាញជាស្រេច។ \\[3pt]
 &  &  &
\ditem ផ្ដាំផ្ញើសុវត្ថិភាពចរាចរណ៍ ជួយការងារឪពុកម្ដាយ។ &
ណែនាំប្រអប់ផ្ទៀងខ្លួនឯង ៣ ចំណុច។ \\ \hline

\rowcolor{stepgray}
\multicolumn{2}{|r|}{\textbf{សរុបរយៈពេលបង្រៀនពេញ}} & \textbf{៥០ នាទី} &
\multicolumn{2}{L{11.50cm}|}{ត្រូវគ្នា ១០០\% នឹង \en{LS\_Plan\_SAS.tex} និង \en{Printable\_Materials\_SAS.tex} (ទំព័រ ១--៦)} \\ \hline
\end{longtable}

\vspace{2pt}

% -------------------- II. បញ្ជីត្រៀមសម្ភារៈ & ការគ្រប់គ្រងស្ថានការណ៍ --------------------
\begin{minipage}[t]{0.46\textwidth}
  \textbf{II.} {\muol បញ្ជីត្រៀមសម្ភារៈមុនចូលថ្នាក់}\\[3pt]
  \begin{tabular}{|c|L{7.0cm}|}
    \hline
    \rowcolor{headgray}
    \textbf{\ding{51}} & \multicolumn{1}{c|}{\textbf{សម្ភារៈត្រូវរៀបចំទុកជាមុន}} \\ \hline
    \chkbox & កញ្ចប់ចង្កឹះ/បន្ទះរឹង $4\text{cm}, 6\text{cm}$ មូលម្ជុលគន្លឹះ \& ក្រដាសស្តើងផ្ទៀងត្រួត \\ \hline
    \chkbox & រ៉ាប៉ទ័រ និងបន្ទាត់ក្រិតសម្រាប់ចាក់សោ $\angle A=60^\circ$ ($BC\approx 5.29\text{cm}$) \\ \hline
    \chkbox & កាតស៊ើបអង្កេតសញ្ញាគន្លឹះ ៣ ជំហាន (ស-ខ្មៅ) \& ផ្ទាំងធនាគារហេតុផល \cnum{1}--\cnum{5} \\ \hline
    \chkbox & Slide Animation ៖ ចាប់កំហុស, ស្រះទឹក, ប៉ោល $C_1,C_2$, មេដៃ ៤ កាត ស-ខ្មៅ \\ \hline
  \end{tabular}
\end{minipage}%
\hfill
\begin{minipage}[t]{0.51\textwidth}
  \textbf{III.} {\muol កំហុសភ័ន្តច្រឡំ \enbf{(}តាមការស្រាវជ្រាវ\enbf{) \&} ដំណោះស្រាយ}\\[3pt]
  \begin{tabular}{|L{8.6cm}|}
    \hline
    \rowcolor{headgray}
    \multicolumn{1}{|c|}{\textbf{ឧបសគ្គសិស្ស (Empirical Misconceptions) \& វិធានការទប់ស្កាត់}} \\ \hline
    \textbf{១. អន្ទាក់ ជ.ជ.ម (SSA) \& អគតិរូបបង្វិល-បញ្ច្រាស៖} ឱ្យយោលប៉ោល $4\text{cm}$ ផ្ទាល់ និងអនុវត្ត \textit{គោលការណ៍កំពូលរួម} លើរូបបង្វិល $140^\circ$ / ត្រួតគ្នា។ \\ \hline
    \textbf{២. រំលងជ្រុងរួម $OM=OM$, មុំទល់កំពូល ឬស្មានគ្មានសញ្ញា៖} តម្រូវ \textbf{Pause-Ask} និយាយមាត់មុនគូស ចង្អុលអូស + គូស $\parallel$ (\cnum{3}) ហាមគូស $AB=A'B'$។ \\ \hline
    \textbf{៣. ច្រឡំ $=$ vs $\cong$ ឬលំដាប់កំពូលខុស៖} បិទធនាគារ \cnum{1}--\cnum{5} ផ្ទៀង \textbf{Pause-Review} លំដាប់ $O\to A\to B$ ទល់ $O\to A'\to B'$។ \\ \hline
    \textbf{៤. យឺតពេលនៅ E2 (១១ នាទី)៖} បង្រួមពិសោធន៍ ២ មក ៣ នាទី និងកាត់មេដៃ E5 មក ២ នាទី ដើម្បីរក្សា ៥០ នាទីពេញ។ \\ \hline
  \end{tabular}
\end{minipage}

\end{document}
